% xdvi -paper 8x6in -xoffset 2in -yoffset 0in zimpl-pittsburgh.dvi 
% dvips -O 1in,-1in zimpl-pittsburgh.dvi 
% ps2pdf -DAutoRotatePages=/None zimpl-pittsburgh.ps
% $Id: zimpl-pittsburgh.tex,v 1.2 2007/01/02 10:54:31 bzfkocht Exp $
\documentclass[semcolor,landscaspe,slidsonly]{seminar}
\renewcommand{\rmdefault}{pmnj}
\renewcommand{\ttdefault}{pcr}
%\renewcommand{\sfdefault}{phv}
\renewcommand{\sfdefault}{pmy}
\usepackage[T1]{fontenc}
\usepackage{textcomp}
\usepackage[small,euler-digits]{eulervm}
\usepackage[latin1]{inputenc}
\usepackage{latexsym}
\usepackage{keyval}
\usepackage{amsmath}
\usepackage{amssymb}
\usepackage{amsfonts}
\usepackage{amsthm}
\usepackage{dcolumn}
\usepackage{stmaryrd}%\lightning
\usepackage[dvips,%
bookmarks,pdffitwindow,pdfcenterwindow=true,pdfstartview=Fit]{hyperref}
\usepackage{graphicx}
\usepackage{xspace}
\usepackage{color}
\usepackage{booktabs}
\usepackage{listings}
\usepackage{multicol}
\usepackage{pstricks,pst-node,pst-tree,pst-grad,pst-char,pstcol}
%
%\usepackage{verbatim}
%
%\usepackage{euler}
%\hoffset0.0in
%\voffset0.0in
%\vsize297mm
%\hsize210mm
\special{papersize=8in,6in}
%\special{! TeXDict begin /landplus90{true}store end }
\renewcommand{\paperwidth}{8in}
\renewcommand{\paperheight}{6in}
\setlength{\slidewidth}{7.6in}
\setlength{\slideheight}{5.6in}
%\setlength{\slideframewidth}{5mm}
%\setlength{\slideframesep}{5mm}
\slidesmag{4}
\extraslideheight{10in}
%\psset{xunit=\semcm,yunit=\semcm}
%\renewcommand{\slidetopmargin}{12mm}
%\renewcommand{\slidebottommargin}{100mm}
%\renewcommand{\slideleftmargin}{5mm}
%\renewcommand{\sliderightmargin}{5mm}
%
\newrgbcolor{lightbluegreen}{.5 .75 .9}
\newrgbcolor{orange}{1 .65 0}
%
\graphicspath{{slides/}}
%
\newslideframe{vortrag}{%
%{\boxput{\rput{30}{\Huge DRAFT}}
%
\psframe[linestyle=none,fillstyle=gradient,gradangle=90,%
 gradbegin=white,gradend=lightbluegreen,gradmidpoint=1]%
 (0,0)(7,5){#1}
% (-11.2,-8)(18.5,9.2)
% \rput[lb](-9.2,8.58){\includegraphics[height=11mm]{ziblogo2.eps}}
% %\rput[lb](-7.2,9.26){\includegraphics[height=7.5mm]{athene.eps}}
% %\rput[lb](-5.4,9.6){\includegraphics[height=5mm]{atesio.eps}}
% \rput[lb](6,9.25){\includegraphics[height=7.9mm]{istlogo2.eps}}
% %\psset{fillstyle=solid,shadow=true,shadowangle=0,shadowsize=1pt,shadowcolor=lila}%\pscharpath[fillcolor=lila]{Momentum}
% \rput[lb](8.5,9.7){\sc\Large\bfseries\color{lila}Momentum}
% %\pswedge[fillstyle=solid,fillcolor=lightblue](17,7){1.6}{75}{130}
% \rput[lb](15.5,7.5){\includegraphics[height=15mm]{kreisel3.eps}}
% %\psframe[fillstyle=none,fillcolor=black,linestyle=solid](-11.1,-9)(18.4,9)%
% \psframe[linestyle=none,fillstyle=solid,fillcolor=white](-11.2,7)(18.5,7.2)
% \rput[lb](-9.2,-7.5){#1}
}
%
%\slideframe{vortrag}
\slideframe{none}
\renewcommand{\labelitemi}{$\blacktriangleright$}
%
%\graphicspath{{../figure/}}
%
\definecolor{seagreen}{rgb}{0.18,0.74,0.56}
\definecolor{darkgreen}{rgb}{0.0,0.35,0.00}
\definecolor{navyblue}{rgb}{0.0,0.0,0.5}
\definecolor{steelblue}{rgb}{0.27,0.51,0.71}
\definecolor{siennabrown}{rgb}{0.63,0.32,0.18}
\definecolor{firebrickred}{rgb}{0.69,0.13,0.13}
\definecolor{gray60}{rgb}{0.6, 0.6, 0.6}
%
%
\lstloadlanguages{C}
\lstdefinelanguage{mps}{%
   keywords={NAME,ROWS,COLUMNS,RHS,BOUNDS,ENDATA},%
   sensitive,%
   keywordstyle=\color{navyblue},%
}[keywords]%
%
\lstdefinelanguage{YACC}{%
   keywords={\%token,\%type,\%left,\%right,\%union},%
   sensitive,%
   singlecomment={/*}{*/},%
   stringizer=[b]',%
   keywordstyle=\color{navyblue},%
   commentstyle=\color{darkgreen},%
   stringstyle=\color{steelblue}%
}[keywords,comments,strings]%
%
\lstdefinelanguage{zimpl}{%
   keywords={set,var,param,minimize,maximize,subto,defnumb},%
   ndkeywords={read,as,comment,binary,integer,real,sum,forall,do,in,proj,vif,vabs,and,or,then,else,end},%
   sensitive,%
   showstringspaces=false,%
   morecomment=[l]\#,
   morestring=[b]",%
   keywordstyle=\color{red},%
   ndkeywordstyle=\color{navyblue},%
   commentstyle=\color{darkgreen},%
   stringstyle=\color{steelblue}%
}[keywords,comments,strings]%
%
\lstdefinestyle{myc}{%
   basicstyle=\sffamily\scriptsize,%
   numberstyle=\sffamily\tiny\color{siennabrown},stepnumber=1,%
   keywordstyle=\color{navyblue},%
   commentstyle=\color{darkgreen},%
   stringstyle=\color{steelblue},%
   directivestyle=\color{firebrickred}%
}
%
\newcommand{\ie}{{i.\,e.,}\xspace}%
\newcommand{\eg}{{e.\,g.}\xspace}%
\newcommand{\BB}{\ensuremath{\{0,1\}}}%
\newcommand{\RR}{\ensuremath{\mathbb{R}}}%
\newcommand{\NN}{\ensuremath{\mathbb{N}}}%
\newcommand{\NNZ}{\ensuremath{\mathbb{N}_0}}%
\newcommand{\ZZ}{\ensuremath{\mathbb{Z}}\xspace}
\newcommand{\QQ}{\ensuremath{\mathbb{Q}}}
\newcommand{\RRP}{\ensuremath{\mathbb{R}_+}}%
\newcommand{\NP}{{\ensuremath{\cal NP}}}
\newcommand{\LP}{LP\xspace}
\newcommand{\IP}{IP\xspace}
\newcommand{\MIP}{MIP\xspace}
\newcommand{\miplib}{miplib\xspace}
\newcommand{\coin}{{\sc coin}\xspace}
\newcommand{\xpress}{XPress\xspace}
\newcommand{\osl}{{\sc osl}\xspace}
\newcommand{\abacus}{{abacus}\xspace}
\newcommand{\concert}{{\sc concert}\xspace}
\newcommand{\mathematica}{{\sc mathematica}\xspace}
\newcommand{\matlab}{{\sc matlab}\xspace}
\newcommand{\zimpl}{{\sc Zimpl}\xspace}
\newcommand{\soplex}{{\sc SoPlex}\xspace}
\newcommand{\perplex}{{\sc perPlex}\xspace}
\newcommand{\porta}{{\sc porta}\xspace}
\newcommand{\sip}{{\sc sip}\xspace}
\newcommand{\gmp}{{\sc gmp}\xspace}
\newcommand{\lp}{{\sc lp}\xspace}
\newcommand{\ip}{{\sc ip}\xspace}
\newcommand{\gnu}{{\sc gnu}\xspace}
\newcommand{\mip}{{\sc mip}\xspace}
\newcommand{\mps}{{\sc mps}\xspace}
\newcommand{\lpf}{{\sc lp}\xspace}
\newcommand{\ibm}{{\sc ibm}\xspace}
\newcommand{\aircom}{{\sc Aircom}\xspace}
\newcommand{\tornado}{{\sc Tornado}\xspace}
\newcommand{\atoll}{{\sc Atoll}\xspace}
\newcommand{\zpl}{{\sc zpl}\xspace}
\newcommand{\hum}{{\sc hum}\xspace}
\newcommand{\ampl}{{\sc ampl}\xspace}
\newcommand{\gams}{{\sc gams}\xspace}
\newcommand{\ilog}{{\sc Ilog}\xspace}
\newcommand{\cplex}{{\sc cplex}\xspace}
\newcommand{\scip}{{\sc scip}\xspace}
\newcommand{\zib}{{\sc zib}\xspace}
\newcommand{\atesio}{{\sf Atesio}\xspace}
\newcommand{\excel}{{\sc excel}\xspace}
\newcommand{\netview}{{\sc NetView}\xspace}
\newcommand{\gmt}{{\sc gmt}\xspace}
\newcommand{\ieee}{{\sc ieee}\xspace}
\newcommand{\code}[1]{{\tt #1}\xspace}
\newcommand{\Order}[1]{\ensuremath{{\cal O}(#1)}}
%
\newcommand{\eplus}{{e$\cdot$plus}\xspace}
\newcommand{\dfn}{{\sc dfn}\xspace}
\newcommand{\tkat}{{\sc Telekom Austria}\xspace}
\newcommand{\bts}{{\sc bts}\xspace}
\newcommand{\bsc}{{\sc bsc}\xspace}
\newcommand{\msc}{{\sc msc}\xspace}
\newcommand{\vlr}{{\sc vlr}\xspace}
\newcommand{\gsm}{{\sc gsm}\xspace}
\newcommand{\poi}{{\sc poi}\xspace}
%
\newcommand{\pentium}{{\sc pentium}\xspace}
\newcommand{\umts}{{\sc umts}\xspace}
\newcommand{\fdd}{{\sc fdd}\xspace}
\newcommand{\www}{{\sc www}\xspace}
\newcommand{\momentum}{{\sc Momentum}\xspace}
\newcommand{\kantenna}{{\sc k742212}\xspace}
\newcommand{\unit}[1]{\,{#1}\xspace} % \sf\small 
\newcommand{\unitx}[1]{#1} % \sf\small 
\newcommand{\smallunit}[1]{\,{\sf\scriptsize #1}\xspace}
\newcommand{\db}{\unit{dB}}
\newcommand{\dbm}{\unit{dBm}}
\newcommand{\degree}{\ensuremath{^\circ}\xspace}%
\newcommand{\wcdma}{{\sc wcdma}\xspace}
\newcommand{\cir}{{\sc cir}\xspace}%
\newcommand{\sites}{\ensuremath{\mathcal{S}}}%
\newcommand{\siteOfInstallation}[1]{\ensuremath{\sigma({#1})}}%
\newcommand{\maxInstallations}[1]{\ensuremath{\Upsilon^{\max}_{#1}}}%
\newcommand{\minInstallations}[1]{\ensuremath{\Upsilon^{\min}_{#1}}}%
\newcommand{\site}[1]{\ensuremath{s_{#1}}}%
\newcommand{\installations}{\ensuremath{\mathcal{I}}}%
\newcommand{\attenuationU}[2]{\ensuremath{\gamma^\uparrow_{#1#2}}}%
\newcommand{\attenuationD}[2]{\ensuremath{\gamma^\downarrow_{#1#2}}}%
%\newcommand{\attenuation}[2]{\ensuremath{\gamma_{{#1},{#2}}}}%
\newcommand{\mobiles}[1]{\ensuremath{\mathcal{M}_{#1}}}%
\newcommand{\mobilesU}[1]{\ensuremath{\mathcal{M}^\uparrow_{#1}}}%
\newcommand{\mobilesD}[1]{\ensuremath{\mathcal{M}^\downarrow_{#1}}}%
\newcommand{\snapshots}{\ensuremath{\mathcal{D}}}%
\newcommand{\totalPowerD}[1]{\ensuremath{\bar{p}^{\downarrow}_{#1}}}%
\newcommand{\totalPowerU}[1]{\ensuremath{\bar{p}^{\uparrow}_{#1}}}%
\newcommand{\powerU}[1]{\ensuremath{p^{\uparrow}_{#1}}}%
\newcommand{\powerD}[2]{\ensuremath{p^{\downarrow}_{#1#2}}}%
\newcommand{\powerP}[1]{\ensuremath{\hat p^{\downarrow}_{#1}}}%
\newcommand{\powerC}[1]{\ensuremath{\check p^{\downarrow}_{#1}}}%
\newcommand{\noise}[1]{\ensuremath{\eta_{#1}}}%
\newcommand{\activityU}[1]{\ensuremath{\alpha^{\uparrow}_{#1}}}%
\newcommand{\activityD}[1]{\ensuremath{\alpha^{\downarrow}_{#1}}}%
\newcommand{\CIRtargetU}[1]{\ensuremath{\mu^{\uparrow}_{#1}}}%
\newcommand{\CIRtargetD}[1]{\ensuremath{\mu^{\downarrow}_{#1}}}%
\newcommand{\CIRtargetP}[1]{\ensuremath{\hat\mu^{\downarrow}_{#1}}}%
\newcommand{\orthogo}[2]{\ensuremath{\omega_{#2}}}%
\newcommand{\invorthogo}[2]{\ensuremath{\bar{\omega}_{#1#2}}}%
\newcommand{\interferenceD}[2]{\ensuremath{\phi({#1},{#2})}}%
\newcommand{\interferenceP}[2]{\ensuremath{\hat\phi({#1},{#2})}}%
\newcommand{\interferenceC}[2]{\ensuremath{\check\phi({#1},{#2})}}%
\newcommand{\loadU}[1]{{l^\uparrow_{#1}}}
\newcommand{\loadD}[1]{{l^\downarrow_{#1}}}
\newcommand{\served}[2]{\ensuremath{x_{#1#2}}}%
\newcommand{\powerLimitUmax}[1]{\ensuremath{\Pi^{\max \uparrow}_{#1}}}%
\newcommand{\powerLimitUmin}[1]{\ensuremath{\Pi^{\min \uparrow}_{#1}}}%
\newcommand{\powerLimitDmaxTotal}[1]{\ensuremath{\Pi^{\max\downarrow}_{#1}}}%
\newcommand{\powerLimitDminLink}[1]{%
   \ensuremath{\dot\Pi^{\min\downarrow}_{#1}}}%
\newcommand{\powerLimitDmaxLink}[1]{%
   \ensuremath{\dot\Pi^{\max\downarrow}_{#1}}}%
\newcommand{\powerLimitDminPilot}[1]{%
   \ensuremath{\hat\Pi^{\min\downarrow}_{#1}}}%
\newcommand{\powerLimitDmaxPilot}[1]{%
   \ensuremath{\hat\Pi^{\max\downarrow}_{#1}}}%
\newcommand{\pilotRscpMin}[1]{\ensuremath{\hat{\pi}^{\min\downarrow}_{#1}}}%
\newcommand{\snapshotOfMobile}[1]{\ensuremath{\rho({#1})}}%
\newcommand{\commonPilotPower}{\ensuremath{\hat{p}^{\downarrow\mbox{c}}}}%
\newcommand{\installation}[1]{\ensuremath{z_{#1}}}%
\newcommand{\rscp}{{\sc RSCP}\xspace}
\newcommand{\CCHfactor}[1]{\ensuremath{\nu_{#1}}}
\newcommand{\pcpich}{pilot\xspace}
\newcommand{\rrm}{{\sc rrm}\xspace}%
\newcommand{\hata}{{\sc cost231-hata}\xspace}
%
\newcommand{\TreeParBox}[2]{\Tr{\psframebox{\parbox{#1}{#2}}}}
\newcommand{\psorgabox}[3]{%
\psframebox[linecolor=blue,linewidth=.1]{\begin{minipage}[t][#1][t]{#2}{\begin{center}\sffamily\scriptsize\vfill%
{#3}\\\vfill\end{center}}\end{minipage}}}
%
\newcommand{\psorgaboxred}[3]{%
\psframebox[linecolor=red,linewidth=.1]{\begin{minipage}[t][#1][t]{#2}{\begin{center}\sffamily\scriptsize\vfill%
{#3}\\\vfill\end{center}}\end{minipage}}}
%
\newcommand{\psorgaboxgray}[3]{%
\psframebox[linecolor=gray,linewidth=.1]{\begin{minipage}[t][#1][t]{#2}{\begin{center}\sffamily\scriptsize\vfill%
{#3}\\\vfill\end{center}}\end{minipage}}}
%
\newcommand{\psorgaboxyel}[3]{%
\psframebox[linecolor=orange,linewidth=.1]{\begin{minipage}[t][#1][t]{#2}{\begin{center}\sffamily\scriptsize\vfill%
{#3}\\\vfill\end{center}}\end{minipage}}}
%
\newcommand{\psorgaboxblue}[3]{%
\psframebox[linecolor=navyblue,linewidth=.1]{\begin{minipage}[t][#1][t]{#2}{\begin{center}\sffamily\scriptsize\vfill%
{#3}\\\vfill\end{center}}\end{minipage}}}
%
\newcommand{\ram}{\textsc{ram}\xspace}
\newcommand{\HVSten}{\textsc{hv}\xspace}
\newcommand{\POI}{\textsc{poi}\xspace}
\newcommand{\VVSten}{\textsc{vv}\xspace}
\newcommand{\UVSten}{\textsc{uv}\xspace}
\newcommand{\OVSten}{\textsc{ov}\xspace}
\newcommand{\NVSten}{\textsc{nv}\xspace}
\newcommand{\HV}{\ensuremath{H}}
\newcommand{\PV}{\ensuremath{P}}
\newcommand{\VV}{\ensuremath{V}}
\newcommand{\UV}{\ensuremath{U}}
\newcommand{\AVH}{\ensuremath{A_{V\!H}}}
\newcommand{\AVP}{\ensuremath{A_{V\!P}}}
\newcommand{\AUV}{\ensuremath{A_{UV}}}
\newcommand{\AAA}{\ensuremath{A}}
% !!!!! ???
\newcommand{\xuv}{\ensuremath{x}}
\newcommand{\xvh}{\ensuremath{x}}
\newcommand{\xvp}{\ensuremath{x}}
%\newcommand{\yvp}{\ensuremath{\bar{y}}}
%\newcommand{\yLvh}{\ensuremath{\bar{y}}}
%\newcommand{\yUvh}{\ensuremath{\dot{y}}}
%\newcommand{\zv}{\ensuremath{z}}
%\newcommand{\zh}{\ensuremath{\dot{z}}}
%\newcommand{\zp}{\ensuremath{\ddot{z}}}
\newcommand{\yvp}{\ensuremath{f}}
\newcommand{\yLvh}{\ensuremath{f^d}}
\newcommand{\yUvh}{\ensuremath{f}}
\newcommand{\zv}{\ensuremath{y}}
\newcommand{\zh}{\ensuremath{y}}
\newcommand{\zp}{\ensuremath{y}}
%
\newcommand{\cxuv}{\ensuremath{c}}
\newcommand{\cxvh}{\ensuremath{c}}
\newcommand{\cxvp}{\ensuremath{c}}
\newcommand{\cyvh}{\ensuremath{\bar{c}}}
\newcommand{\cyvp}{\ensuremath{\bar{c}}}
\newcommand{\czv}{\ensuremath{c}}
\newcommand{\czh}{\ensuremath{c}}
\newcommand{\czp}{\ensuremath{c}}
%\newcommand{\dline}{\ensuremath{\bar{\delta}}}
%\newcommand{\duser}{\ensuremath{\dot{\delta}}}
\newcommand{\dline}{\ensuremath{\delta^d}}
\newcommand{\duser}{\ensuremath{\dot{\delta}}}
%\newcommand{\maxvvline}{\ensuremath{\bar{\mu}}}
%\newcommand{\maxvvuser}{\ensuremath{\dot{\mu}}}
%\newcommand{\maxhvline}{\ensuremath{\bar{\nu}}}
%\newcommand{\maxhvuser}{\ensuremath{\dot{\nu}}}
\newcommand{\maxvvline}{\ensuremath{\rho^d_v}}
\newcommand{\maxvvuser}{\ensuremath{\dot{\mu}}}
\newcommand{\maxhvline}{\ensuremath{\rho^d_h}}
\newcommand{\maxhvuser}{\ensuremath{\dot{\nu}}}
%
\newcommand{\poitrans}{\ensuremath{\tau}}
%
\newcommand{\Code}[1]{\texttt{#1}\xspace}
\newcommand{\pots}{\textsc{pots}\xspace}
\newcommand{\isdn}{\textsc{isdn}\xspace}
%
% dfn -----
%
\newcommand{\gwin}{\textsf{\small G-WiN}\xspace}
\newcommand{\bwin}{\textsf{\small B-WiN}\xspace}
%
\newcommand{\setR}{\ensuremath{R}\xspace}
\newcommand{\setU}{\ensuremath{U}\xspace}
\newcommand{\setV}{\ensuremath{W}\xspace}
\newcommand{\setW}{\ensuremath{V}\xspace}
\newcommand{\setA}[1]{\ensuremath{A_{#1}}\xspace}
\newcommand{\setS}[1]{\ensuremath{S_{#1}}\xspace}
\newcommand{\Auv}{\ensuremath{\setA{\setU\setV}}\xspace}
\newcommand{\Avw}{\ensuremath{\setA{\setV\setW}}\xspace}
\newcommand{\Sv}{\ensuremath{S_{\vv}}\xspace}
\newcommand{\Sw}{\ensuremath{S_{\ww}}\xspace}
\newcommand{\uu}{\ensuremath{u}\space}
\newcommand{\vv}{\ensuremath{w}\space}
\newcommand{\ww}{\ensuremath{v}\space}
\newcommand{\rr}{\ensuremath{r}\space}
\newcommand{\si}{\ensuremath{s}\space}
%
%\newcommand{\x}[1]{\ensuremath{x_{#1}}\xspace}
%\newcommand{\y}[1]{\ensuremath{f^{\rr}_{#1}}\xspace}
%\newcommand{\z}[1]{\ensuremath{z^{\si}_{#1}}\xspace}
%\newcommand{\xuv}{\ensuremath{x_{\uu\vv}}\xspace}
\newcommand{\xvw}{\ensuremath{x_{\vv\ww}}\xspace}
\newcommand{\yrvw}{\ensuremath{f^{\rr}_{\vv\ww}}\xspace}
\newcommand{\yrw}{\ensuremath{f^{\rr}_{\ww}}\xspace}
\newcommand{\zs}[1]{\ensuremath{z^{\si}_{#1}}\xspace}
\newcommand{\zsv}{\ensuremath{z^{\si}_{\vv}}\xspace}
\newcommand{\zsw}{\ensuremath{z^{\si}_{\ww}}\xspace}
%
\newcommand{\dru}{\ensuremath{\delta^{\rr}_{\uu}}\xspace}
\newcommand{\crs}[1]{\ensuremath{\kappa^{\rr\si}_{#1}}\xspace}
\newcommand{\crsv}{\ensuremath{\kappa^{\rr\si}_{\vv}}\xspace}
\newcommand{\crsw}{\ensuremath{\kappa^{\rr\si}_{\ww}}\xspace}
\newcommand{\rvmax}{\ensuremath{\lambda^{\rr}_{\vv}}\xspace}
%
\newcommand{\obj}{\ensuremath{\tau}\xspace}
%
\newcommand{\stp}{{\sc stp}\xspace}
\newcommand{\stpp}{{\sc stpp}\xspace}
\newcommand{\tsp}{{\sc tsp}\xspace}
\newcommand{\vlsi}{{\sc vlsi}\xspace}
\newcommand{\nets}{{\cal N}\xspace}
\newcommand{\fa}{\ensuremath{\text{for all }}}
\newcommand{\with}{\ensuremath{\text{\ with\ }}}
\newcommand{\tal}{{\sc tal}\xspace}
\newcommand{\mol}{{\sc mol}\xspace}
%
% implementation
%
\newcommand{\discnet}{{\sc discnet}\xspace}
\newcommand{\lex}{{\sc lex}\xspace}
\newcommand{\flex}{{\sc flex}\xspace}
\newcommand{\bison}{{\sc bison}\xspace}
\newcommand{\yacc}{{\sc yacc}\xspace}
\newcommand{\unix}{{\sc unix}\xspace}
\newcommand{\LALR}{{\sc lalr}\xspace}
\newcommand{\posix}{{\sc posix}\xspace}
\newcommand{\netlib}{{\sc netlib}\xspace}
%\newcommand{\}{{\sc }\xspace}
%
\newcommand{\cpu}{{\sc cpu}\xspace}
\newcommand{\zsets}{\ensuremath{{\cal\check{S}}_n}\xspace}
\newcommand{\zset} [1]{\ensuremath{\check{#1}_n}\xspace}
\newcommand{\ntuple}{\ensuremath{{\cal T}_n}\xspace}
\newcommand{\abs}{\ensuremath{\operatorname{abs}}}
\newcommand{\sgn}{\ensuremath{\operatorname{sgn}}}
\newcommand{\loc}{{\sc loc}\xspace}
\newcommand{\lint}{{\sc lint}\xspace}
\newcommand{\flexelint}{{\sc flexelint}\xspace}
\newcommand{\valgrind}{{\sc valgrind}\xspace}
%
\newcommand{\PO}{\ensuremath{P}\xspace} %{\textsl{P}\xspace}

\newpsobject{showgrid}{psgrid}{subgriddiv=1,griddots=10,gridlabels=1pt}
\newcommand{\slidesubtitlea}{Subtitlea}
\newcommand{\slidesubtitleb}{Subtitleb}
\newcommand{\slidetitle}[1]{%
\centerslidesfalse%
\sffamily%
\psset{xunit=1\semin,yunit=1\semin}
\begin{pspicture}(0,0)(8,0.5)%\showgrid
\rput[bl](7.3,-0.08){\includegraphics[width=0.4\semin]{zib-logo}}%
\rput[bl](6.8,0.15){\includegraphics[width=0.45\semin]{matheon-logo}}%
\rput[bl](0,0.2){\color{blue}#1}%
\rput[bl](4,0.4){\tiny\sc\slidesubtitlea}%
\rput[bl](4,0.2){\tiny\sc\slidesubtitleb}%
\psline[linewidth=0.5pt](0,0.1)(7.7,0.1)%
\end{pspicture}%
}
%
\onlyslides{}
\pagestyle{empty}
\begin{document}
%
\setlength{\arraycolsep}{2pt}
\renewcommand{\slidesubtitlea}{INFORMS 2006, Pittsburgh, 7. Nov 2006}%
\renewcommand{\slidesubtitleb}{Thorsten Koch, Zuse Institute Berlin}%
%
\begin{slide}
\slidetitle{{\small Open Source Modeling Tools}}

\medskip
\Large
\begin{center}
{{\LARGE\bfseries\textsc{Zimpl}}

\large\vspace{0.4\semin}
{\color{red}A Solver Independend Open Source\\Modeling Language}}

\large\vspace{0.7\semin}
{Thorsten Koch, ZIB, \textsc{Matheon}}
\end{center}
\end{slide}

\begin{slide}
{\it %Try to find the real tense of the talk you are hearing: Was it done,
 %is it being done, or is something to be done? \ldots 

 CONGRAM (CONtractor GRAMmer):

\begin{center}
 defined by the {\color{blue}imperfect past},\\ 
 the {\color{blue}insufficient present}, and the\\
 {\color{blue}absolutely perfect future}.
\end{center}
}
\begin{flushright}
--- fortune(6)
\end{flushright}
\end{slide}
\begin{slide}
\slidetitle{}
\vskip2cm
\begin{center}
\LARGE Imperfect Past
\end{center}
\end{slide}


%-------------------------------------------------------------------------
% What is a MPL 
%-------------------------------------------------------------------------
\begin{slide}
\slidetitle{Algebraic modeling language}
\small
\zimpl is an algebraic modeling language 
for mixed integer programming problems.

Modeling languages allow to formulate a
mathematical programming problem in 
a way similar to the original mathematical formulation and 
automatically generate input for a corresponding solver.

\begin{itemize}
\item Problems are represented in a declarative way
\item Separation between problem definition and solution process
\item Separation between problem structure and data.
\end{itemize}

Algebraic modeling languages typically operate on
contraints consisting of variables indexed by sets connected to
parameters by algebraic expressions.
%% \normalsize Uses for MPL's

%% \small
%% \begin{itemize}
%% \item Teaching modeling\\
%%       easy to learn and use
%% \item Practical use\\
%%       rapid prototyping
%% \item Research\\
%%       results are easier to reproduce
%% \end{itemize}

\end{slide}

\begin{slide}
%\slidetitle{Modeling language}
\lstset{language=zimpl,%
basicstyle=\sffamily\scriptsize,%
numberstyle=\sffamily\tiny\color{siennabrown},stepnumber=0}
\begin{lstlisting}{}
# This model solves Sudoku puzzels.
param p := 3;

set L := { 1 .. p*p };
set M := { 1 .. p};

var x[L*L*L] binary;

# file: row col value
set F := { read "sudoku.dat" as "<1n,2n,3n>" comment "#" };

subto dots: forall <i,j> in L*L do sum <k> in L : x[i,j,k] == 1;
subto rows: forall <j,k> in L*L do sum <i> in L : x[i,j,k] == 1;
subto cols: forall <i,k> in L*L do sum <j> in L : x[i,j,k] == 1;
subto sqrs: forall <m,n,k> in M*M*L do 
   sum <i,j> in M*M : x[(m-1)*p+i,(n-1)*p+j,k] == 1;

# Fix the fixed values
subto fixed: forall <i,j,k> in F do x[i,j,k] == 1;

\end{lstlisting}
\end{slide}
%-------------------------------------------------------------------------
% Beispiel 1
%-------------------------------------------------------------------------
\begin{slide}
\slidetitle{Modeling language}
\small
For $a,b,c\in\ZZ$, 
$a\in [-15,15], b\in [-10,20], c\in [-20,10]$,

maximize $5 a + 3 b + c$
subject to:
\begin{equation*}
\begin{array}{rl}
\text{If}&(a \neq b \text{ and }(|a - b| = 3 c \text{ or } c - a \geq 0))\\
 \text{then}&a + b + c \geq 7\\
 \text{else}& a + b\leq 1
\end{array}
\end{equation*}

\bigskip
{\color{blue}can be modeled as an IP, but this very laborious\ldots}
\end{slide}

%-------------------------------------------------------------------------
% Beispiel 2 
%-------------------------------------------------------------------------
\begin{slide}
\slidetitle{Modeling language}
\small
%\frqq{}Erweiterte\flqq{} Funktionen:
For $a,b,c\in\ZZ$, 
$a\in [-15,15], b\in [-10,20], c\in [-20,10]$,

maximize $5 a + 3 b + c$ 
subject to:
\begin{equation*}
\begin{array}{rl}
\text{If}&(a \neq b \text{ and }(|a - b| = 3 c \text{ or } c - a \geq 0))\\
 \text{then}&a + b + c \geq 7\\
 \text{else}& a + b\leq 1
\end{array}
\end{equation*}

\smallskip
\lstset{language=zimpl,%
basicstyle=\sffamily\footnotesize,%
numberstyle=\sffamily\tiny\color{siennabrown},stepnumber=0}
\begin{lstlisting}[frame=t]{}
var a integer >= -15 <= 15;
var b integer >= -10 <= 20;
var c integer >= -20 <= 10;
maximize obj: 5 * a + 3 * b + c;
subto c1: 
   vif (a!=b and (vabs(a-b) == 3*c or c-a >= 0)) 
      then a + b + c >= 7 
      else a + b <= 1 end;
\end{lstlisting}
\end{slide}

%-------------------------------------------------------------------------
% Ergebnis 
%-------------------------------------------------------------------------
%\centerslidestrue
\begin{slide}
\begin{minipage}{0.5\textwidth}
\tiny
\begin{verbatim}
- xp0  + bp2                      <= 0
- b    + a       - xp0  + bm3     <= 0
- xp0  + 25 bp2                   >= 0
- b    + a       - xp0  + 35 bm3  >= 0
- bp2  - bm3     + re4             = 0
- b    + a       + xm6  - xp5      = 0
- xp5  + 25 bp7                   >= 0
+ xm6  + 35 bp7                   <= 35
- xm6  - xp5     + re8             = 0
- xp9  + bp11                     <= 0
- 3 c  + re8     - xp9  + bm12    <= 0
- xp9  + 95 bp11                  >= 0
- 3 c  + re8     - xp9  + 30 bm12 >= 0
+ bp11 + bm12    + re13            = 1
- xp14 + bp16                     <= 0
+ c    - a       - xp14 + bm17    <= 0
- xp14 + 25 bp16                  >= 0
+ c    - a       - xp14 + 35 bm17 >= 0
+ bp16 + bm17                     <= 1
+ re13 - re19                     <= 0
- bm17 - re19                     <= -1
+ re13 - bm17    - re19           >= -1
+ re4  - re20                     >= 0
+ re19 - re20                     >= 0
+ re4  + re19    - re20           <= 1
+ c    + b       + a    - 52 re20 >= -45
+ b    + a              - 34 re20 <= 1
\end{verbatim}
\end{minipage}\hfill
\begin{minipage}{0.4\textwidth}
\sffamily
\scriptsize
IP after presolve: 

27 rows, 20 columns, 77 non-zeros.

\medskip
{\color{blue}Optimal solution:}

\begin{tabular}{lr}
%\toprule
{\color{red}c}      &             {\color{red} 2}\\
{\color{red}b}      &             {\color{red}20}\\
{\color{red}a}      &             {\color{red}14}\\
xm1   &              6\\
bm3   &              1\\
re4   &              1\\
xm6   &              6\\
re8   &              6\\
re13  &             1\\
xm15  &            12\\
bm17  &             1\\
re19  &             1\\
re20  &             1\\
%\bottomrule
\end{tabular}

\end{minipage}
\end{slide}


\begin{slide}
\slidetitle{Why another MPL}

Modeling languages have been developed since the 80's.

Examples: AIMMS, AMPL, GAMS, LINGO, LPL, MOSEL, MPL, \ldots

Why yet another modeling language?
\end{slide}

\begin{slide}
\slidetitle{Why another MPL}

Modeling languages have been developed since the 80's.

Examples: AIMMS, AMPL, GAMS, LINGO, LPL, MOSEL, MPL, \ldots

Why yet another modeling language?

\bigskip
{\it\color{blue} Software is like sex. It's better when it's free.$^*$}
\begin{flushright}
--- Linux Torwalds
\end{flushright}
\end{slide}

\begin{slide}
\slidetitle{Why another MPL}

Modeling languages have been developed since the 80's.

Examples: AIMMS, AMPL, GAMS, LINGO, LPL, MOSEL, MPL, \ldots

Why yet another modeling language?

\bigskip
{\it\color{blue} Software is like sex. It's better when it's free.$^*$}
\begin{flushright}
--- Linux Torwalds
\end{flushright}

\vfill
{\scriptsize $^*$ Maybe it is not so sophisticated, 
but there is more love and dedication involved.}
\end{slide}
%-------------------------------------------------------------------------
% What is a MPL 
%-------------------------------------------------------------------------
\begin{slide}
\slidetitle{}
\vskip2cm
\begin{center}
\LARGE Insufficent Present
\end{center}
\end{slide}

\begin{slide}
\slidetitle{{\large\zimpl}}
%{\large\zimpl} is
\begin{itemize}
\setlength{\itemsep}{0pt}
\item is freely available with C source code under the GPL.
\item is highly portable (*NIX, Windows, MacOS-X).
\item is solver independend.
\item is available as a library.
\item can be used standalone or linked to a solver.
\item is pretty stable. 
\item has been used in several lectures and industry projects.
\item can count correctly.
\end{itemize}
\end{slide}

\begin{slide}
\slidetitle{Freely available}

In contrast to most other modeling languages, {\color{blue}\zimpl is 
freely available in source code under GPL.}

This is very convienent for teaching, as all students can use Zimpl
on their laptops and at home. 

\zimpl is highly portable and {\color{blue}compiles under Linux, Windows, MacOS-X,
Solaris, Tru64, IRIX, HP-UX, \ldots}

Only dependend on the GMP (GNU Multi Precision Library).

Download at {\color{red}http://zimpl.zib.de}, binaries and a user
guide are available.

\end{slide}

%-------------------------------------------------------------------------
% Solver independend
%-------------------------------------------------------------------------
\begin{slide}
\slidetitle{Solver independend}

\zimpl used standalone generates either MPS or LP files, which can be
processed by most solvers.

\zimpl can write SOS and ORD files for CPLEX.

It is strait forward and easy to add further formats\\
(as far as \zimpl is concerned ;-)

%If \zimpl is used as a library and linked to a solver 
%Currently, all features are supported by all solvers, but this might
%change. 

\end{slide}

%-------------------------------------------------------------------------
% Zimpl Statistik
%-------------------------------------------------------------------------
\begin{slide}
\slidetitle{Stable}
\centering
{\sffamily\scriptsize
\begin{tabular}{>{\raggedleft}p{2cm}rrrrrr}\toprule
\multicolumn{1}{l}{\emph{Zimpl statistics}}&\color{blue}LOC&\color{blue}Stmt.&\color{blue}Calls&\color{blue}CC&\color{blue}Ass.&\color{blue}Cover\\
\midrule
684 functions, total & 11369 & 7520 & 4398 & 1972 & 1255 & 86\%\\
$\varnothing$ per function& 16.6 & 11.0 & 6.4 & 2.9 & 1.8 &\\
$\varnothing$ statements per& 0.7 && 1.7 & 3.8 & 6 &\\
%\midrule
%\multicolumn{7}{l}{\emph{PerPlex statistics for comparison}}\\
%\midrule
%%163 functions, total & 3847 & 2617 & 1292 & 836 & 315\\
%$\varnothing$ per function  & 23.6 & 16.1 & 7.9 & 5.1 & 1.9 &\\
%$\varnothing$ statements per& 0.7  && 2.0 & 3.1 & 8 &\\
\bottomrule
\end{tabular}
}

\vfill
\includegraphics[width=\textwidth]{coverage}

\end{slide}


\begin{slide}
\slidetitle{Use as a library}

\lstset{language=C,style=myc}%
%basicstyle=\sffamily\footnotesize,%
%numberstyle=\sffamily\tiny\color{siennabrown},stepnumber=0}
\begin{lstlisting}{}
Bool zpl_read_with_args(int argc, char** argv);

Con* xlp_addcon(const char* name, ConType type,
   const Numb* lhs, const Numb* rhs, unsigned int flags);
Var* xlp_addvar(const char* name, VarClass varclass,
   const Bound* lower, const Bound* upper, 
   const Numb* priority, const Numb* startval);
void xlp_addtonzo(Var* var, Con* von, const Numb* numb);

Sos*   xlp_addsos(const char* name, 
   SosType type, const Numb* priority);
void   xlp_addtocost(Var* var, const Numb* cost);
Bound* xlp_getlower(const Var* var);
Bound* xlp_getupper(const Var* var);
void   xlp_setdir(Bool minimize);
\end{lstlisting}

\end{slide}

\begin{slide}
\tiny
\begin{verbatim}
SCIP version 0.90a [precision: 8 byte] [mode: optimized] [LP solver: SOPLEX 1.3.0]
Copyright (c) 2002-2006 Konrad-Zuse-Zentrum fuer Informationstechnik Berlin (ZIB)

SCIP> r decide2.zpl

****************************************************
* Zuse Institute Mathematical Programming Language *
* Release 2.05  Copyright (C)2006 by Thorsten Koch *
****************************************************

original problem has 105 variables (105 bin, 0 int, 0 impl, 0 cont) and 28 constraints
SCIP> opt

presolving: ...

 time | node  | left  |LP iter| mem |frac |strbr|  dualbound   | primalbound  |  gap   
r 0.0s|     1 |     0 |    40 | 285k|  21 |   0 | 1.965773e+02 | 1.226000e+03 | 523.67%
R 0.0s|     1 |     0 |    40 | 286k|  21 |   0 | 1.965773e+02 | 1.208000e+03 | 514.52%
...
* 0.0s|     1 |     0 |   109 | 294k|   - |  13 | 2.010000e+02 | 2.010000e+02 |   0.00%

SCIP Status        : problem is solved [optimal solution found]
Solving Time (sec) : 0.04
Solving Nodes      : 1
Primal Bound       : +2.01000000000000e+02 (9 solutions)
Dual Bound         : +2.01000000000000e+02
Gap                : 0.00 %
\end{verbatim}
\end{slide}
\begin{slide}
\tiny
\begin{verbatim}
> bin/cbc-generic decide2.zpl
Default solver is clp
****************************************************
* Zuse Institute Mathematical Programming Language *
* Release 2.05  Copyright (C)2006 by Thorsten Koch *
****************************************************

Cgl0004I processed model has 28 rows, 100 columns (100 integer) and 618 elements
processed model has 28 rows and 100 columns
Cbc0012I Integer solution of 201 found by heuristic after 8 iterations
and 0 nodes (0.02 seconds)
Cbc0031I 3 added rows had average density of 93
Cbc0013I At root node, 3 cuts changed objective from 196.577 to 201 in 2 passes
Cbc0014I Cut generator 0 (Probing) - 0 row cuts (0 active), 0 column
cuts  in 0.008001 seconds - new frequency is 10
Cbc0014I Cut generator 1 (Gomory) - 15 row cuts (3 active), 0 column
cuts  in 1.00289e-18 seconds - new frequency is 1
Cbc0014I Cut generator 2 (Knapsack) - 0 row cuts (0 active), 0 column
cuts  in 1.00289e-18 seconds - new frequency is 100
Cbc0014I Cut generator 3 (Clique) - 0 row cuts (0 active), 0 column
cuts  in 1.00289e-18 seconds - new frequency is 100
Cbc0014I Cut generator 4 (MixedintegerRounding) - 0 row cuts (0 active),
0 column cuts  in 0.004 seconds - new frequency is 100
Cbc0014I Cut generator 5 (FlowCover) - 0 row cuts (0 active), 0 column
cuts  in -2.73083e-18 seconds - new frequency is 100
Cbc0001I Search completed - best objective 201, took 8 iterations and 0
nodes (0.02 seconds)
Result 201 iterations 8 nodes 0 took 0.024001 seconds - total 0.056003
\end{verbatim}
\end{slide}

\begin{slide}
\slidetitle{Uses}
{\large\color{red}Zimpl has (and is) being used for:}
\begin{itemize}
\setlength{\itemsep}{0pt}%
\item {\color{blue}Location planning for DFN, e-plus and Telekom Austria,}
\item {\color{blue}UMTS location planning}
\item {\color{blue}Generating 3D-Steiner tree packing problem,}
\item Track auctioning,
\item Protein folding,
\item \ldots,
\item {\color{blue}several lectures, e.g. at TU Berlin, Uni Bayreuth,
  courses in Berlin and Beijing,} and
\item works for Models with more than a million variables.
\end{itemize}
\end{slide}

\begin{slide}
\slidetitle{}
\vskip2cm
\begin{center}
\LARGE Absolutly Perfect Future
\end{center}
\end{slide}
%-------------------------------------------------------------------------
% Differences 1
%-------------------------------------------------------------------------
\begin{slide}
\slidetitle{What is not on the list?}

\begin{itemize}
\item Integrated windows environment\\ 
(possible if somebody has time for it)

\item Not the boss, but just a reader/translator

\item Not procedural

\item Trying to beat GAMS or AMPL on features.
%Only a subset of the possibilities of, e.g. AMPL.
\end{itemize}
\end{slide}

%-------------------------------------------------------------------------
% Differences 2
%-------------------------------------------------------------------------
\begin{slide}
\slidetitle{Next on the list}

\begin{itemize}
\setlength{\itemsep}{0pt}%
\item Integrating with more solvers:
COIN CBC (works), lp\_solve (experimental), 
SQopt\_ex (most interesting: compete rational environment)
\item Reduce memory requirements for large models.
\item More extended functions
\item more attributes, e.g. semi continous variables,
      implicit integer variables
\item support quadratic objective functions
\item specific constraints
\item general constraint interface (passthrough to solver)
\item complete rational preprocessing
\end{itemize}  

\end{slide}

\begin{slide}
\slidetitle{}
\vskip2cm
\begin{center}
\LARGE Do not wait for it!
\end{center}
\end{slide}

%-------------------------------------------------------------------------
% Differences 2
%-------------------------------------------------------------------------
\begin{slide}
\slidetitle{Summary}
\vskip1cm
\begin{center}
\bfseries{{\Large\zimpl} 

is a small

\color{blue}free} and {\color{blue}usable} 

modeling language for MIP 

{\color{red}you can use today.}
\end{center}
\end{slide}
%-------------------------------------------------------------------------
% Thank you
%-------------------------------------------------------------------------
\begin{slide}
\slidetitle{}

\vfill
\Huge\bfseries\centering Thank you!

\vfill
\end{slide}

%-------------------------------------------------------------------------
% Thank you
%-------------------------------------------------------------------------
\begin{slide}
\slidetitle{}
\Huge\bfseries

\vfill
\centering Questions?

\vfill
\centering Comments?

\vfill
\end{slide}
\end{document}
%-------------------------------------------------------------------------
% Modellierungssprachen machen alles leichter
%-------------------------------------------------------------------------
\centerslidesfalse
\begin{slide}
\sffamily
%{\color{red}Die Verwendung einer Modellierungssprache erleichtert die 
%Ver�ffentlichung und erh�ht die Reproduzierbarkeit:}
%
\bfseries
{\color{blue}Vorteile durch Verwendung von Modellierungssprachen}

\smallskip
\hrule

\bigskip
\begin{itemize}

\item {\color{red}Modellierungssprachen erleichtern das (schnelle) Experimentieren
      mit Ideen und Modellen.}% Das ist Hilfreich bei Industrieprojekten.

\item Die Reproduzierbarkeit der Ergebnisse wird erh�ht.

\item Die �bersetzung eines Modells von einer Modellierungssprache in eine andere ist
      meist einfach.

\item Gr��ere Unabh�ngigkeit vom verwendeten L�ser.

%\item Erweiterte Funktionen sind problematisch.

\item {\color{blue}Es ist m�glich, die Umsetzung erweiterter Funktionen
      l�serabh�ngig zu machen.}
\end{itemize}
\end{slide}

%-------------------------------------------------------------------------
% Modellierungssprachen erleichtern das Experimentieren
%-------------------------------------------------------------------------
\centerslidestrue
\begin{slide}
\sffamily
%Mathematik/Anwendung/Lehre\\
%
%\vfill
\Large
\begin{center}
{\color{blue}Vorteile durch Verwendung von Modellierungssprachen:}

Modellierungssprachen erleichtern das Experimentieren
mit Ideen und Modellen.
\end{center}
%\vfill
\end{slide}

%-------------------------------------------------------------------------
% Zimpl im Einsatz
%-------------------------------------------------------------------------
\begin{slide}
\sffamily\bfseries
{\large\color{red}Zimpl wurde und wird verwendet f�r:}
\begin{itemize}
\setlength{\itemsep}{0pt}%
\item {\color{blue}Standortplanungsprobleme beim DFN, e-plus\\ und der Telekom Austria,}
\item {\color{blue}UMTS Planung,}
\item {\color{blue}3D-Steinerbaumpackungsprobleme,}
\item Trassenb�rse,
\item Proteinfaltung,
\item \ldots,
\item verschiedene Vorlesungen,
\item und Modelle mit mehr als 10 Millionen Variablen.
\end{itemize}

{\small\color{firebrickred}Website {\tt\color{darkgreen} http://www.zib.de/koch/zimpl}}

\end{slide}


\end{document}

%-------------------------------------------------------------------------
% Einfach zu lernen und zu benutzen
%-------------------------------------------------------------------------
\begin{slide}
\sffamily
%Mathematik/Anwendung/Lehre\\
%
%\vfill
\Large
\begin{center}
Wie sind die\\ erweiterten Funktionen modelliert?
\end{center}
%\vfill
\end{slide}

%-------------------------------------------------------------------------
% Implementierung der erweiterten Funktionen
%-------------------------------------------------------------------------
\centerslidesfalse
\begin{slide}
\sffamily
{\color{blue}Modellierung der erweiterten Funktionen}

\medskip
\hrule

\bigskip
Gegeben eine beschr�nkte ganzzahlige Variable 

\begin{center}
{\color{red}$l_x\leq x\leq u_x$} mit $l_x,x,u_x\in\ZZ$
\end{center}
\begin{equation*}
\begin{array}{rcccl}
&&\color{blue}x^+ - x^-&=&x\\
b^+&\le&x^+&\le&\max(0,u_x) b^+\\
b^-&\le&x^-&\le&|\min(0,l_x)|b^-\\
&&\color{blue}b^+ + b^- & \le & 1\\
&&b^+,b^-&\in&\BB
\end{array}\label{eqn:conditionals}
\end{equation*}

\hrule
\smallskip
\small
\begin{minipage}{0.45\textwidth}
\begin{tabular}{ll}
\color{blue}$x^+$&positiver Anteil von $x$\\
\color{blue}$x^-$&negativer Anteil von $x$\\
\end{tabular}
\end{minipage}\hfill
\begin{minipage}{0.45\textwidth}
\begin{tabular}{ll}
\color{blue}$b^+$&Indikator f�r $x>0$\\
\color{blue}$b^-$&Indikator f�r $x<0$\\
\end{tabular}
\end{minipage}
\end{slide}

%-------------------------------------------------------------------------
% Implementierung 2 
%-------------------------------------------------------------------------
\begin{slide}
\sffamily
{\color{blue}Modellierung der erweiterten Funktionen}

\medskip
\hrule

\smallskip
%Die folgenden Funktionen und Relationen k�nnen wie folgt modelliert werden:

\begin{minipage}{0.4\textwidth}
$$
\begin{array}{lcl}
\abs(x)   & = & x^+ + x^-\\
\sgn(x)   & = & b^+ - b^-\\
\min(v,w) & = & w - x^-\\
\max(v,w) & = & x^+ + w\\
\end{array}
$$
\end{minipage}\hfill
\begin{minipage}{0.5\textwidth}
$$
\begin{array}{lcl}
v\neq w   & \Leftrightarrow & b^+ + b^- = 1\\
v =   w   & \Leftrightarrow & b^+ + b^- = 0\\
v \leq  w & \Leftrightarrow & b^+ = 0\\
v <     w & \Leftrightarrow & b^- = 1\\
v \geq  w & \Leftrightarrow & b^- = 0\\
v >     w & \Leftrightarrow & b^+ = 1 
\end{array}
$$
\end{minipage}

\medskip
wobei 
$x=v-w$ f�r $v,w\in\ZZ$ mit $l_x=l_v-u_w$ und $u_x=u_v-l_w$ gilt, falls
zwei Operanden beteiligt sind.

\end{slide}

%-------------------------------------------------------------------------
% vif
%-------------------------------------------------------------------------
\begin{slide}
\sffamily
{\color{blue}Modellierung von vif mittels big-M}

\medskip
\hrule

\bigskip
Gegeben eine Variable $r\in\{0,1\}$ und eine Bedingung 
$$\sum a_i x_i\leq b$$ 
mit $-\infty\leq l_i \leq x_i \leq u_i\leq\infty$ und
$$U=\sum_{a_i<0} a_i l_{i} + \sum_{a_i>0} a_i u_{i}<\infty$$

Die Bedingung braucht nur dann erf�llt sein, wenn $r=1$:

\begin{equation*}
\sum a_i x_i + (U-b)r\le U\label{eqn:cond-exec}
\end{equation*}

\end{slide}

%-------------------------------------------------------------------------
% Boolesche Operatoren
%-------------------------------------------------------------------------
\begin{slide}
\sffamily
{\color{blue}Modellierung von booleschen Operatoren}

\medskip
\hrule

\bigskip
\begin{minipage}{0.48\textwidth}
{\color{red}$a$ = $b$ and $c$}
\begin{equation*}
\begin{array}{rcl}
a - b&\le&0\\
a - c&\le&0\\
a - b - c&\ge&-1%\\
%a&=&b \mbox{ and } c\\
\end{array}
\label{eqn:modell-and}
\end{equation*}

{\color{red}$a$ = $b$ or $c$}
\begin{equation*}
\begin{array}{rcl}
a - b&\ge&0\\
a - c&\ge&0\\
a - b - c&\le&0%\\
%a&=&b \mbox{ or } c\\
\end{array}
\end{equation*}
\end{minipage}\hfill
\begin{minipage}{0.48\textwidth}
{\color{red}$a$ = not $b$}
\begin{equation*}
\begin{array}{rcl}
a + b&=&1%\\
%a&=&\neg b\\
\end{array}
\end{equation*}

{\color{red}$a$ = $b$ xor $c$}
\begin{equation*}
\begin{array}{rcl}
a - b - c&\le&0\\
a - b + c&\ge&0\\
a + b - c&\ge&0\\
a + b + c&\le&2%\\
%a&=&b \mbox{ xor } c\\
\end{array}
\end{equation*}
% {\footnotesize Alternative 1}
% \begin{equation*}
% a+b+c=2d
% \end{equation*}
% {\footnotesize Alternative 2}
% \begin{equation*}
% \begin{array}{rcl}
% x-b+c&=&0\\
% x-m+n&=&0\\
% a-m-n&=&0
% \end{array}
% \end{equation*}
%Note, that the polyhedron of the linear relaxation of this alternative
%formulation has non integral vertices, \eg, $a=b=0$, $c=1$, and $d=0.5$.

\end{minipage}
\end{slide}
